% Generated by python -m scripts.render_exact_catalogue; exact Q(sqrt d) identities.
\paragraph{2O 的无理二面角表。}
下表的 $\theta[T]\in(0,\pi/2)$ 满足 $\tan^2\theta[T]=T$；
$\cos\theta<0$ 的二面角取 $\pi-\theta[T]$。
{\small\begin{longtable}{@{}ll@{\hspace{3em}}ll@{}}\toprule
$T=\tan^2\theta$ & $\theta$ & $T=\tan^2\theta$ & $\theta$\\\midrule\endhead
$\tfrac{1}{8}$ & $\tfrac{1}{2}\pi-2\beta$ & $\tfrac{1}{2}$ & $\beta$\\
$2$ & $\tfrac{1}{2}\pi-\beta$ & $8$ & $2\beta$\\
\bottomrule\end{longtable}}
\paragraph{2O 的证书（3 条精确恒等式）。}
{\small\begin{longtable}{@{}rll@{}}\toprule 类型 & 关系 & \\\midrule\endhead
互余 & $\theta[\tfrac{1}{8}]+\theta[8]=\tfrac{1}{2}\pi$ & \\
互余 & $\theta[\tfrac{1}{2}]+\theta[2]=\tfrac{1}{2}\pi$ & \\
差角 & $\theta[8]-2\theta[\tfrac{1}{2}]=0$ & \\
\bottomrule\end{longtable}}
\paragraph{2I 的无理二面角表。}
下表的 $\theta[T]\in(0,\pi/2)$ 满足 $\tan^2\theta[T]=T$；
$\cos\theta<0$ 的二面角取 $\pi-\theta[T]$。
{\small\begin{longtable}{@{}ll@{\hspace{3em}}ll@{}}\toprule
$T=\tan^2\theta$ & $\theta$ & $T=\tan^2\theta$ & $\theta$\\\midrule\endhead
$\tfrac{1}{5}$ & $B$ & $\tfrac{1}{4}$ & $A$\\
$\tfrac{3}{5}$ & $\tfrac{1}{2}C$ & $\tfrac{4}{5}$ & $\tfrac{1}{2}\pi-2B$\\
$\tfrac{5}{4}$ & $2B$ & $\tfrac{5}{3}$ & $\tfrac{1}{2}\pi-\tfrac{1}{2}C$\\
$4$ & $\tfrac{1}{2}\pi-A$ & $5$ & $\tfrac{1}{2}\pi-B$\\
$15$ & $C$ & $\tfrac{63}{121}-\tfrac{24}{121}\sqrt5$ & $-\tfrac{1}{3}\pi+C$\\
$\tfrac{63}{121}+\tfrac{24}{121}\sqrt5$ & $\tfrac{2}{3}\pi-C$ & $\tfrac{7}{8}-\tfrac{3}{8}\sqrt5$ & $-\tfrac{1}{2}A+B$\\
$\tfrac{7}{8}+\tfrac{3}{8}\sqrt5$ & $\tfrac{1}{2}\pi-\tfrac{1}{2}A-B$ & $\tfrac{3}{2}-\tfrac{1}{2}\sqrt5$ & $\tfrac{1}{4}\pi-\tfrac{1}{2}A$\\
$\tfrac{3}{2}+\tfrac{1}{2}\sqrt5$ & $\tfrac{1}{4}\pi+\tfrac{1}{2}A$ & $3-\tfrac{4}{3}\sqrt5$ & $-\tfrac{1}{6}\pi+\tfrac{1}{2}C$\\
$3+\tfrac{4}{3}\sqrt5$ & $\tfrac{1}{6}\pi+\tfrac{1}{2}C$ & $\tfrac{7}{2}-\tfrac{3}{2}\sqrt5$ & $\tfrac{1}{4}\pi-B$\\
$\tfrac{7}{2}+\tfrac{3}{2}\sqrt5$ & $\tfrac{1}{4}\pi+B$ & $9-4\sqrt5$ & $\tfrac{1}{2}A$\\
$9+4\sqrt5$ & $\tfrac{1}{2}\pi-\tfrac{1}{2}A$ & $14-6\sqrt5$ & $\tfrac{1}{2}A+B$\\
$14+6\sqrt5$ & $\tfrac{1}{2}\pi+\tfrac{1}{2}A-B$ & $27-12\sqrt5$ & $\tfrac{1}{3}\pi-\tfrac{1}{2}C$\\
$27+12\sqrt5$ & $\tfrac{2}{3}\pi-\tfrac{1}{2}C$ & &\\
\bottomrule\end{longtable}}
\paragraph{2I 的证书（22 条精确恒等式）。}
{\small\begin{longtable}{@{}rll@{}}\toprule 类型 & 关系 & \\\midrule\endhead
互余 & $\theta[\tfrac{1}{5}]+\theta[5]=\tfrac{1}{2}\pi$ & \\
互余 & $\theta[\tfrac{1}{4}]+\theta[4]=\tfrac{1}{2}\pi$ & \\
互余 & $\theta[\tfrac{3}{5}]+\theta[\tfrac{5}{3}]=\tfrac{1}{2}\pi$ & \\
互余 & $\theta[\tfrac{4}{5}]+\theta[\tfrac{5}{4}]=\tfrac{1}{2}\pi$ & \\
互余 & $\theta[\tfrac{7}{8}-\tfrac{3}{8}\sqrt5]+\theta[14+6\sqrt5]=\tfrac{1}{2}\pi$ & \\
互余 & $\theta[\tfrac{7}{8}+\tfrac{3}{8}\sqrt5]+\theta[14-6\sqrt5]=\tfrac{1}{2}\pi$ & \\
互余 & $\theta[\tfrac{3}{2}-\tfrac{1}{2}\sqrt5]+\theta[\tfrac{3}{2}+\tfrac{1}{2}\sqrt5]=\tfrac{1}{2}\pi$ & \\
互余 & $\theta[3-\tfrac{4}{3}\sqrt5]+\theta[27+12\sqrt5]=\tfrac{1}{2}\pi$ & \\
互余 & $\theta[3+\tfrac{4}{3}\sqrt5]+\theta[27-12\sqrt5]=\tfrac{1}{2}\pi$ & \\
互余 & $\theta[\tfrac{7}{2}-\tfrac{3}{2}\sqrt5]+\theta[\tfrac{7}{2}+\tfrac{3}{2}\sqrt5]=\tfrac{1}{2}\pi$ & \\
互余 & $\theta[9-4\sqrt5]+\theta[9+4\sqrt5]=\tfrac{1}{2}\pi$ & \\
差角 & $\theta[\tfrac{5}{4}]-2\theta[\tfrac{1}{5}]=0$ & \\
差角 & $\theta[\tfrac{7}{2}+\tfrac{3}{2}\sqrt5]-\theta[\tfrac{1}{5}]=\tfrac{1}{4}\pi$ & \\
差角 & $\theta[\tfrac{1}{4}]-2\theta[9-4\sqrt5]=0$ & \\
差角 & $\theta[\tfrac{63}{121}-\tfrac{24}{121}\sqrt5]-2\theta[3-\tfrac{4}{3}\sqrt5]=0$ & \\
差角 & $\theta[15]-\theta[\tfrac{63}{121}-\tfrac{24}{121}\sqrt5]=\tfrac{1}{3}\pi$ & \\
差角 & $\theta[\tfrac{63}{121}+\tfrac{24}{121}\sqrt5]-2\theta[27-12\sqrt5]=0$ & \\
差角 & $\theta[\tfrac{3}{5}]-\theta[3-\tfrac{4}{3}\sqrt5]=\tfrac{1}{6}\pi$ & \\
差角 & $\theta[3+\tfrac{4}{3}\sqrt5]-\theta[\tfrac{3}{5}]=\tfrac{1}{6}\pi$ & \\
差角 & $\theta[4]-2\theta[\tfrac{3}{2}-\tfrac{1}{2}\sqrt5]=0$ & \\
差角 & $\theta[\tfrac{1}{5}]-\theta[\tfrac{7}{8}-\tfrac{3}{8}\sqrt5]-\theta[9-4\sqrt5]=0$ & \\
差角 & $\theta[9+4\sqrt5]-\theta[\tfrac{1}{5}]-\theta[\tfrac{7}{8}+\tfrac{3}{8}\sqrt5]=0$ & \\
\bottomrule\end{longtable}}
